\documentclass[11pt]{article}
\usepackage[T2A]{fontenc}
\usepackage[utf8]{inputenc}
\usepackage[russian]{babel}
\usepackage{amsmath,amssymb}
\usepackage[margin=2cm]{geometry}

\title{Формулы КР: расчёт информационной цепочки}
\author{Папин А.В., ИУ5Ц-31М}
\date{2026}

\begin{document}
\maketitle

\section*{Цепочка}
кадр $\rightarrow$ камера $\rightarrow$ детектор $\rightarrow$
представление $\rightarrow$ оператор $\rightarrow$ решение

Итоговая точность:
\[
A = \prod_{i=1}^{n} a_i
\]

Вклад звена в потерю:
\[
\ell_i = \frac{1 - a_i}{\sum_j (1 - a_j)}
\]

\section*{Калибровка уверенности}
\[
\mathrm{ECE} = \sum_{m=1}^{M} \frac{|B_m|}{n}\,\bigl|\,\mathrm{acc}(B_m)
- \mathrm{conf}(B_m)\,\bigr|
\]
\[
BS = \frac{1}{N}\sum_{i=1}^{N} (f_i - o_i)^2
\]

\section*{Теория обнаружения сигналов}
\[
d' = z(H) - z(F), \qquad c = -\tfrac{1}{2}\bigl(z(H) + z(F)\bigr)
\]

\section*{Восприятие и время}
\[
RT = a + b\,\log_2(n+1)
\]
\[
MT = a + b\,\log_2(A/W + 1)
\]

\section*{Различимость}
\[
\Delta I / I = k
\]
\[
\theta = 2\arctan\!\left(\frac{h}{2d}\right)
\]
\[
\Delta E_{00}, \qquad \frac{L_1 + 0.05}{L_2 + 0.05}
\]

\section*{Свежесть и нагрузка}
\[
\Delta(t) = t - u(t), \qquad 7 \pm 2, \quad 4 \pm 1
\]
\[
\mathrm{TLX} = \sum_k w_k R_k, \qquad A \propto S \cdot E_f \cdot E_x \cdot V
\]

\end{document}
